\documentclass[11pt]{article}
\usepackage{amsmath,amssymb,amsthm,hyperref}
\newtheorem{theorem}{Theorem}
\newtheorem{lemma}[theorem]{Lemma}
\title{An Exact Integer-Fiber Singular-Jet Compiler}
\author{}
\date{}
\begin{document}
\maketitle

\paragraph{Status.}
The results below are complete written finite constructions in the original
attention model. They are not Lean-checked. They do not establish an
asymptotic improvement for arbitrary Boolean functions. The final section
states the resource accounting that prevents treating the construction as
an asymptotic solution. An exact small-instance audit is supplied separately.

\begin{theorem}[A whole-cube zero representation]
Let $w_1,\ldots,w_n$ be positive integers, put
$T(x)=\sum_i w_i x_i$ and $R=1+\sum_i w_i$, and let $H>R$.
There is a sum of $H$ nonzero original admissible constant-numerator heads
which vanishes at every point of $\{0,1\}^n$. All denominator slopes are
strictly negative. All coefficients can be rational.
\end{theorem}
\begin{proof}
Put $d_t=1-t/(4R)$ for $0\le t<R$, choose distinct rational poles
$0<\lambda_h<1/4$, and set
\[
 P(z)=\prod_{t=0}^{R-1}(z-d_t),\qquad
 Q(z)=\prod_{h=1}^H(z-\lambda_h),\qquad
 c_h=\frac{P(\lambda_h)}{Q'(\lambda_h)}.
\]
Partial fractions give
\[
 \sum_{h=1}^H\frac{c_h}{z-\lambda_h}=\frac{P(z)}{Q(z)}.
 \tag{1}
\]
Substitute $z=d_{T(x)}$. The right side vanishes. Each $c_h$ is
nonzero because the poles are distinct and disjoint from the roots of $P$.
The denominators are
\[
 1-\lambda_h-\sum_i\frac{w_i}{4R}x_i>\frac12.
 \tag{2}
\]
Their original-atom realization is included in the proof below.
\end{proof}

\begin{theorem}[Independent polynomial readouts on integer fibers]
Keep $T,R,d_t$ as above, let $K\ge0$, and prescribe, independently for
each $0\le t<R$, an affine function $L_t:\mathbb R^n\to\mathbb R$
and a real polynomial $p_t$ of degree at most $K$. Put
\[
 H=(2K+1)R-K.
 \tag{3}
\]
There is a family $F_\varepsilon$ of sums of $H$ original admissible
constant-numerator heads, defined for all sufficiently small
$\varepsilon>0$, such that, for every Boolean $x$ with $T(x)=t$,
\[
 \lim_{\varepsilon\downarrow0}\varepsilon^{-K}F_\varepsilon(x)
   =\frac{p_t(L_t(x))}{Q(d_t)}.
 \tag{4}
\]
Here $Q(d_t)>0$ is independent of the chosen readouts. Consequently, if
every prescribed $p_{T(x)}(L_{T(x)}(x))$ is nonzero, some such $H$-head
sum realizes all their signs strictly. Rational readouts admit rational
head parameters.
\end{theorem}

\begin{proof}
Choose distinct poles $0<\lambda_h<1/4$, and define $Q$ as above and
$P_0(z)=\prod_{t=0}^{R-1}(z-d_t)$. Interpolate the affine readouts in
the variable $z$: there is a polynomial $L(z,x)$ of degree at most
$R-1$ in $z$, affine in $x$, with $L(d_t,x)=L_t(x)$.
Write $p_t(s)=\sum_{j=0}^K a_{tj}s^j$. For each $j$ choose the unique
polynomial $A_j$ of degree at most $R-1$ with
\[
 A_j(d_t)=\frac{a_{tj}}{P_0'(d_t)^j}.
 \tag{5}
\]
The roots $d_t$ are distinct, so all these interpolation data are defined.
Put
\[
 P_\varepsilon(z)=\sum_{j=0}^K
   \varepsilon^{K-j}A_j(z)P_0(z)^j,
 \qquad
 F_\varepsilon(x)=\sum_{h=1}^H
 \frac{P_\varepsilon(\lambda_h)/Q'(\lambda_h)}
 {d_{T(x)}-\lambda_h+\varepsilon L(\lambda_h,x)}.
 \tag{6}
\]
For $0\le k\le K$ the polynomial
$P_\varepsilon(z)L(z,x)^k$ has degree at most
\[
 (K+1)R-1+K(R-1)=(2K+1)R-K-1=H-1.
 \tag{7}
\]
Thus the partial-fraction identity (1), followed by $k$ derivatives,
gives
\[
 \sum_h\frac{P_\varepsilon(\lambda_h)L(\lambda_h,x)^k}
 {Q'(\lambda_h)(z-\lambda_h)^{k+1}}
 =\frac{(-1)^k}{k!}\frac{d^k}{dz^k}
 \left(\frac{P_\varepsilon(z)L(z,x)^k}{Q(z)}\right).
 \tag{8}
\]
Fix $x$ and let $t=T(x)$. A finite geometric expansion of each
denominator in (6), using the positive fixed gaps $d_t-\lambda_h$,
therefore yields
\[
 F_\varepsilon(x)=
 \sum_{k=0}^K\frac{\varepsilon^k}{k!}
 \left[\frac{d^k}{dz^k}
 \left(\frac{P_\varepsilon(z)L(z,x)^k}{Q(z)}\right)\right]_{z=d_t}
 +O(\varepsilon^{K+1}).
 \tag{9}
\]
The remainder statement follows directly from the finite geometric
remainder formula: $P_\varepsilon(\lambda_h)$ is bounded as
$\varepsilon\to0$, there are finitely many heads, and their denominators
stay bounded away from zero. No interchange of infinite series is used.

Insert the $j$th summand of $P_\varepsilon$ into the $k$th term of (9).
If $j>k$, its derivative vanishes because $P_0(z)^j$ has a zero of
order $j$ at $d_t$. If $j<k$, its power of $\varepsilon$ is
$K-j+k>K$. If $j=k$, the only surviving term of the $k$th derivative
puts all derivatives on $P_0^k$ and equals
\[
 k!\frac{A_k(d_t)P_0'(d_t)^kL_t(x)^k}{Q(d_t)}
 =k!\frac{a_{tk}L_t(x)^k}{Q(d_t)}.
 \tag{10}
\]
Equations (9)--(10) prove (4).

It remains to check the exact original model. For each head, write its
denominator as $b_h(\varepsilon)-\sum_i a_{hi}(\varepsilon)x_i$.
At $\varepsilon=0$,
\[
 b_h(0)=1-\lambda_h>\frac34,\qquad
 a_{hi}(0)=\frac{w_i}{4R}>0,\qquad
 b_h(0)-2\sum_i a_{hi}(0)>\frac14.
 \tag{11}
\]
All these inequalities persist for sufficiently small $\varepsilon>0$,
because the finitely many affine coefficients depend continuously on
$\varepsilon$. Set
\[
 \alpha_h=\frac12,\quad \rho_{hi}=2a_{hi}(\varepsilon),\quad
 \gamma_h=b_h(\varepsilon)-2\sum_i a_{hi}(\varepsilon),\quad
 m_{hi}=0,\quad \delta_h=0,\quad
 \eta_h=\frac{P_\varepsilon(\lambda_h)}{Q'(\lambda_h)}.
 \tag{12}
\]
These are original admissible atom parameters. On Boolean inputs,
$\rho_{hi}\alpha_h^{x_i}=2a_{hi}-a_{hi}x_i$, so their denominator
and numerator are exactly the corresponding summand of (6).
In particular, admissibility holds on the entire cube, including fibers
not used by an application. Since the cube is finite, nonzero limits in
(4) preserve all prescribed strict signs for one common small positive
$\varepsilon$. Rational input data and a sufficiently small rational
$\varepsilon$ give rational parameters throughout.
\end{proof}

\paragraph{Useful collision families and the actual cost.}
For $m$ row bits $u$ and $b$ memory bits $v$, choose
\[
 T(u,v)=\sum_{i=0}^{m-1}2^i u_i+\sum_{j=1}^b a_jv_j,
 \qquad a_j\in\mathbb Z_{>0}.
 \tag{13}
\]
Then $R=2^m+\sum_j a_j$. If $\sum_j a_j=O(2^m)$, the zero
representation and the degree-$K$ compiler cost respectively $O(2^m)$
and $O(K2^m)$ heads. On a fiber, the row bits encode
$t-\sum_j a_jv_j$, whenever this integer is in $[0,2^m-1]$.
Thus arbitrary weights permit binary-carry geometry, rather than only
dependence on memory Hamming weight.

The theorem does not assert that one polynomial of one affine readout
programs arbitrary signs on a fiber. Nor does it identify the independent
readouts with free coefficients in the original final score: their
interpolation enters the denominators in (6), and the cancellation proof
is essential. To obtain the requested logarithmic saving from this
compiler one would need fibers of size substantially greater than
$nK$ that admit all sign patterns as $p_t(L_t(x))$. A generic
polynomial-activation capacity slogan does not supply such a theorem.
Repeated roots alone are therefore not an asymptotic solution.

\begin{theorem}[Capacity obstruction for strict leading readouts]
Let $n,K\ge1$, and let $X\subseteq\{0,1\}^n$. Suppose that every strict
sign pattern on $X$ is realized by $p(L(x))$, with $L$ affine and
$\deg p\le K$. Then $|X|<128nK$.
Consequently, for a fixed partition into $R$ integer fibers, the strict
leading readouts in the preceding compiler cannot represent every Boolean
function with $KR=o(2^n/n)$.
\end{theorem}
\begin{proof}
Write $M=|X|$, $p_*=n+K+2$, and $D=K+1$. At a fixed point $x$,
the expression $\sum_{j=0}^K a_j(\ell_0+\sum_i\ell_i x_i)^j$
is a polynomial of total degree at most $D$ in the $p_*$ parameters
$a_j,\ell_i$. We use Warren's strict sign-pattern estimate: $M\ge p_*$
polynomials of degree at most $D$ in $p_*$ real variables have at most
$(4eDM/p_*)^{p_*}$ nonzero sign patterns. This is the standard form
quoted as Theorem 2.1 in Goldberg and Jerrum,
\emph{Bounding the Vapnik--Chervonenkis dimension of concept classes
parameterized by real numbers},
\url{https://www.cs.ox.ac.uk/people/paul.goldberg/papers/goldbergjerrum.pdf}.
Thus shattering implies, when $M\ge p_*$,
\[
 2^{M/p_*}\le 4eD(M/p_*).
 \tag{14}
\]
For $a\ge1$, the inequality $2^z\le az$ forces
$z<4\log_2(2a)$: $2^z/z$ is increasing for $z\ge2$, while at
$z=4\log_2(2a)$ it is greater than $a$, using
$\log_2(2a)\le2a$. The case $M<p_*$ is smaller still. Hence
\[
 M<4(n+K+2)\log_2(8e(K+1)).
 \tag{15}
\]
If $K\le n$, then $n+K+2\le4n$ and
$\log_2(8e(K+1))<6K$, giving $M<96nK$.
If $n\le K\le2^n$, then $n+K+2\le4K$ and
$\log_2(8e(K+1))<7n$, giving $M<112nK$.
If $K>2^n$, then $M\le2^n<K\le nK$.
These bounds imply the stated uniform constant. If the independent
readout class on all $R$ fibers realizes every global sign pattern,
each fiber must be shattered by this same local family. Summing the
bound over fibers gives $2^n<128nKR$.
\end{proof}

\paragraph{Exact boundary of the obstruction.}
This last theorem concerns the nonzero strict limiting readouts in (4).
It is not a lower bound for the original attention model. A construction
that deliberately makes leading readouts zero and uses further orders,
or that shares a root lattice between many levels at smaller head cost,
is not excluded. Such a construction requires a new cancellation argument;
neither multiplicity nor interpolation alone supplies it.
\end{document}
